%-----------------------------------------------------------------------------------------------------------------------------------------------%
%	The MIT License (MIT)
%
%	Copyright (c) 2015 Jan Küster
%
%	Permission is hereby granted, free of charge, to any person obtaining a copy
%	of this software and associated documentation files (the "Software"), to deal
%	in the Software without restriction, including without limitation the rights
%	to use, copy, modify, merge, publish, distribute, sublicense, and/or sell
%	copies of the Software, and to permit persons to whom the Software is
%	furnished to do so, subject to the following conditions:
%	
%	THE SOFTWARE IS PROVIDED "AS IS", WITHOUT WARRANTY OF ANY KIND, EXPRESS OR
%	IMPLIED, INCLUDING BUT NOT LIMITED TO THE WARRANTIES OF MERCHANTABILITY,
%	FITNESS FOR A PARTICULAR PURPOSE AND NONINFRINGEMENT. IN NO EVENT SHALL THE
%	AUTHORS OR COPYRIGHT HOLDERS BE LIABLE FOR ANY CLAIM, DAMAGES OR OTHER
%	LIABILITY, WHETHER IN AN ACTION OF CONTRACT, TORT OR OTHERWISE, ARISING FROM,
%	OUT OF OR IN CONNECTION WITH THE SOFTWARE OR THE USE OR OTHER DEALINGS IN
%	THE SOFTWARE.
%	
%
%-----------------------------------------------------------------------------------------------------------------------------------------------%

% Modified By Steven Yu 2025
%
%============================================================================%
%
%	DOCUMENT DEFINITION
%
%============================================================================%

%we use article class because we want to fully customize the page and dont use a cv template
\documentclass[10pt,A4]{article}	


%----------------------------------------------------------------------------------------
%	ENCODING
%----------------------------------------------------------------------------------------

%we use utf8 since we want to build from any machine
\usepackage[utf8]{inputenc}		
\usepackage[UTF8]{ctex}
%----------------------------------------------------------------------------------------
%	LOGIC
%----------------------------------------------------------------------------------------

% provides \isempty test
\usepackage{xifthen}

%----------------------------------------------------------------------------------------
%	FONT
%----------------------------------------------------------------------------------------

% some tex-live fonts - choose your own

%\usepackage[defaultsans]{droidsans}
%\usepackage[default]{comfortaa}
%\usepackage{cmbright}
\usepackage[default]{lato}
%\usepackage{fetamont}
%\usepackage[default]{gillius}
%\usepackage[light,math]{iwona}
%\usepackage[thin]{roboto} 

% set font default
\renewcommand*\familydefault{\sfdefault} 	
\usepackage[T1]{fontenc}

% more font size definitions
\usepackage{moresize}		


%----------------------------------------------------------------------------------------
%	PAGE LAYOUT  DEFINITIONS
%----------------------------------------------------------------------------------------

%debug page outer frames
%\usepackage{showframe}			


%define page styles using geometry
\usepackage[a4paper]{geometry}		

% for example, change the margins to 2 inches all round
\geometry{top=1.25cm, bottom=-.6cm, left=1.5cm, right=1.5cm} 	

%use customized header
\usepackage{fancyhdr}				
\pagestyle{fancy}

%less space between header and content
\setlength{\headheight}{-5pt}		


%customize entries left, center and right
%\lhead{}
%\chead{}
%\rhead{}


%indentation is zero
\setlength{\parindent}{0mm}

%----------------------------------------------------------------------------------------
%	TABLE /ARRAY DEFINITIONS
%---------------------------------------------------------------------------------------- 

%for layouting tables
\usepackage{multicol}			
\usepackage{multirow}

%extended aligning of tabular cells
\usepackage{array}

\newcolumntype{x}[1]{%
>{\raggedleft\hspace{0pt}}p{#1}}%


%----------------------------------------------------------------------------------------
%	GRAPHICS DEFINITIONS
%---------------------------------------------------------------------------------------- 

%for header image
\usepackage{graphicx}

%for floating figures
\usepackage{wrapfig}
\usepackage{float}
%\floatstyle{boxed} 
%\restylefloat{figure}

%for drawing graphics		
\usepackage{tikz}				
\usetikzlibrary{shapes, backgrounds,mindmap, trees}


%----------------------------------------------------------------------------------------
%	Color DEFINITIONS
%---------------------------------------------------------------------------------------- 

\usepackage{color}

%accent color
\definecolor{NJAUGREEN}{RGB}{0,82,68}

%dark background color
\definecolor{bgcol}{RGB}{110,110,110}

%light background / accent color
\definecolor{softcol}{RGB}{238,219,65}

%NJAUBLEU color
\definecolor{NJAUBLEU}{RGB}{4,54,144}


%============================================================================%
%
%
%	DEFINITIONS
%
%
%============================================================================%

%----------------------------------------------------------------------------------------
% 	HEADER
%----------------------------------------------------------------------------------------

% remove top header line
\renewcommand{\headrulewidth}{0pt} 

%remove botttom header line
\renewcommand{\footrulewidth}{0pt}	  	

%remove pagenum
\renewcommand{\thepage}{}	

%remove section num		
\renewcommand{\thesection}{}			

%----------------------------------------------------------------------------------------
% 	ARROW GRAPHICS in Tikz
%----------------------------------------------------------------------------------------

% a six pointed arrow poiting to the left
\newcommand{\tzlarrow}{(0,0) -- (0.2,0) -- (0.3,0.2) -- (0.2,0.4) -- (0,0.4) -- (0.1,0.2) -- cycle;}	

% include the left arrow into a tikz picture
% param1: fill color
%
\newcommand{\larrow}[1]
{\begin{tikzpicture}[scale=0.58]
	 \filldraw[fill=#1!100,draw=#1!100!black]  \tzlarrow
 \end{tikzpicture}
}

% a six pointed arrow poiting to the right
\newcommand{\tzrarrow}{ (0,0.2) -- (0.1,0) -- (0.3,0) -- (0.2,0.2) -- (0.3,0.4) -- (0.1,0.4) -- cycle;}

% include the right arrow into a tikz picture
% param1: fill color
%
\newcommand{\rarrow}
{\begin{tikzpicture}[scale=0.7]
	\filldraw[fill=NJAUGREEN!100,draw=NJAUGREEN!100!black] \tzrarrow
 \end{tikzpicture}
}


\setlength{\headheight}{-5pt}		


%customize entries left, center and right
\lhead{}
\chead{ \small{ \textbf{Phone:} {\textbf{Your Phone Number }}  \textperiodcentered  \textbf{Email:} {\textbf{youremail@gmail.comn }}\textperiodcentered \textbf{GitHub:} \textcolor{NJAUGREEN}{\textbf{https://github.com/seydlitz2185}} }}
\rhead{}


%indentation is zero
\setlength{\parindent}{0mm}
%----------------------------------------------------------------------------------------
%	custom sections
%----------------------------------------------------------------------------------------

% create a coloured box with arrow and title as cv section headline
% param 1: section title
%
\newcommand{\cvsection}[1]
{
	%\begin{center} Center the secton title
	\begin{flushleft}
		\large\textcolor{NJAUGREEN}{\textbf{#1}}
	\end{flushleft}
	%\end{center}
}

%create a coloured arrow with title as cv meta section section
% param 1: meta section title
%
\newcommand{\metasection}[2]
{
%\begin{tabular*}{1\textwidth}{r r}
\footnotesize{#2} \hspace*{\fill} \footnotesize{#1}\\[1pt]
%\end{tabular*}
}

%----------------------------------------------------------------------------------------
%	 CV EVENT
%----------------------------------------------------------------------------------------

% creates a stretched box as cv entry headline followed by some paragraphs about 
% the work you did
% param 1:	event time i.e. 2014 or 2011-2014 etc.
% param 2:	event name (what did you do?)
% param 3:	institution (where did you work / study)
% param 4:	list of paragraphs
%
\newcommand{\cvevent}[4]
{

\begin{tabular*}{1\textwidth}{p{13.6cm}  x{3.9cm}}
	\textbf{#2} - \textcolor{bgcol}{#3} &   \vspace{2.5pt}\textcolor{NJAUGREEN}{#1}
\end{tabular*}

\vspace{-8pt}
\textcolor{softcol}{\hrule}
\vspace{6pt}

	\foreach \desc in {#4}{
		$\cdot$ \desc\\[3pt]
	}
	
\vspace{3pt}
}

% creates a stretched box as 
\newcommand{\cveventmeta}[2]
{
	\mbox{\mystrut \hspace{87pt}\textit{#1}}\\
	#2
}

%----------------------------------------------------------------------------------------
% CUSTOM STRUT FOR EMPTY BOXES
%----------------------------------------- -----------------------------------------------
\newcommand{\mystrut}{\rule[-.3\baselineskip]{0pt}{\baselineskip}}

%----------------------------------------------------------------------------------------
% CUSTOM LOREM IPSUM
%----------------------------------------------------------------------------------------
\newcommand{\lorem}
{Lorem ipsum dolor sit amet, consectetur adipiscing elit. Donec a diam lectus.}



%============================================================================%
%
%
%
%	DOCUMENT CONTENT
%
%
%
%============================================================================%
\begin{document}


%use our custom fancy header definitions
\pagestyle{fancy}	
\vspace{2pt}
\hspace{-0.25\linewidth}{\makebox[1.5\linewidth][c]{\HUGE{\textcolor{NJAUBLEU}{\textbf{姓名}} } \textcolor{NJAUGREEN}{\rule[-1mm]{1mm}{0.9cm}} \HUGE{\textcolor{NJAUBLEU}{\textbf{简历}} } }}


%----------------------------------------------------------------------------------------
%	HEADER IMAGE
%----------------------------------------------------------------------------------------

%\begin{figure}[H]
%\begin{flushright}
%	\includegraphics[trim= 320 130 460 210,clip,width=0.2\linewidth]{myfoto.jpg}	%trimming relative to image size!
%\end{flushright}
%\end{figure}
\begin{figure}[H]
	\begin{flushright}
		\includegraphics[width=0.2\linewidth,angle=0,trim=0 0 0 -100,clip]{linux.jpeg}% Adjust the trim values to fine-tune the position
	\end{flushright}
\end{figure}
	

%---------------------------------------------------------------------------------------
%	TITLE HEADLINE
%----------------------------------------------------------------------------------------
\vspace{-175pt}
\begin{center}
	 \small 出生日期：YYYY-MM-DD |
	\small 年龄：Age |
	 \small 性别：性别 |
	 \small 籍贯：XXXX |
	 \small 联系电话： Your Phone Number |
	\small 邮箱：youremail@gmail.com
\end{center}



\vspace{1pt}


%---------------------------------------------------------------------------------------
%	META SECTION
%----------------------------------------------------------------------------------------
\metasection{}{\textbf{毕业院校：} 你的学校("211"工程，“双一流”一流学科建设高校)}
\metasection{}{\textbf{专业：} 你的专业} 
\metasection{}{\textbf{学位:} 你的学位}
\metasection{}{\textbf{平均学分绩点:} 3.33/4.5 \hspace*{53pt} 算术平均分：85.56}{}
\metasection{}{\textbf{专业课程及成绩:} 数据结构 ：85；  \hspace*{25pt}计算机网络基础：93； \hspace*{25pt}数据库系统：73；\\ \hspace*{58pt} JAVA项目实战：94； \hspace*{5pt}B/S应用系统开发：95； \hspace*{14pt}管理信息系统实践：92；}
\metasection{}{\textbf{考研初试报名专业：}xx大学\hspace*{2pt}计算机学院\hspace*{2pt}085404\hspace*{2pt}计算机技术}
\metasection{}{\textbf{考研初试成绩：}（101）政治：69； \hspace*{52pt}（201）英语（一）：73；\\ \hspace*{52pt}（301）数学（一）：104； \hspace*{27.5pt}（408）计算机专业基础综合：93；\\ \hspace*{58pt}总分：339}

\vspace{-15pt}

\textcolor{softcol}{\hrule}

\cvsection{研究意愿}
\vspace{-5pt}
{结合本科阶段掌握的XXXXXX知识，我希望硕士研究生阶段能进行XXXXX方向的研究。此外，我也将始终保持学习新兴技术、探索前沿知识的热情，愿意积极尝试XXXX领域的研究。}

%---------------------------------------------------------------------------------------
%	SUMMARAY (optional)
%----------------------------------------------------------------------------------------


\vspace{-5pt}
\textcolor{softcol}{\hrule}
%============================================================================%
%
%	CV SECTIONS AND EVENTS (MAIN CONTENT)
%
%============================================================================%

%---------------------------------------------------------------------------------------
%	EXPERIENCE
%----------------------------------------------------------------------------------------
\cvsection{实践经历}
\vspace{-5pt}
\cvevent{YYYY.MM}{你的大学XXXX届本科优秀毕业论文(设计) XX等奖}{Yourl University}{
	{论文题目：你的论文题目 \hspace*{238pt}指导教师：你的导师},
	{内容简介：概括你的毕业设计主要内容。},
	{公示网站：\textcolor{NJAUGREEN}{https://链接}}
}

%
\cvevent{YYYY.MM-YYYY.MM}{XXXX公司}{ aistudio.baidu.com}{
	{内容简介：使用XXXX数据集， 训练XXXX模型，对XXXX进行XXXX研究。\hspace*{40pt}项目导师：你的导师},
	{参与工作：1.数据集预处理：使用Python语言批量压缩图像并调整bbox，生成VOC和coco格式数据集。\\ \hspace*{45pt}2. 编写项目文档：将数据预处理和配置训练参数的过程编写为notebook文档，以供后续活动参考。},
	{文档链接：\textcolor{NJAUGREEN}{https://链接}}
}
\cvevent{YYYY.MM-YYYY.MM}{XXXX 实习}{Nanjing Agricultural University}{
	{实习内容：实习期间主要做了什么。}
}

\vspace{-20pt}
\textcolor{softcol}{\hrule}
%---------------------------------------------------------------------------------------
%	CERTIFICATION SECTION
%--------------------------------------------------------------------------------------
\cvsection{技能特长}
\vspace{-5pt}
\cvevent{}{技能证书}{Skill Certifications}{
	{全国计算机技术与软件专业技术资格（水平）考试（软考）\hspace*{2pt}中级：软件设计师；\hspace*{2pt}初级：程序员},
	{全国计算机等级考试（NCRE）\hspace*{2pt}三级： 数据库技术 \hspace*{2pt}等级：良好},
	{全国大学英语六级考试（CET-6） },
	{普通话水平测试\hspace*{2pt}：一级乙等}
}
\cvevent{}{特长爱好}{Strengths and Hobbies}{
	{特长：熟悉Python自然语言处理、Window Subsystem for Linux、Docker、Git的基本操作；\\\hspace*{5pt}了解基于Vue 3、 Tomcat、 Servelt、MyBatis框架的Java Web系统开发；了解MySQL，Redis数据库基本操作；\\\hspace*{5pt}了解c++基本语法，熟练阅读各种英文文档，具有日常英语交流能力。},
	{爱好：Proxmox VE虚拟机，爬山，跑步。}
}

\vspace{-20pt}
\textcolor{softcol}{\hrule}
\cvsection{自我评价}
\vspace{-5pt}
	{具有较强的自学能力，乐于接触和尝试多种计算机技术。敢于钻研，善于思考，遇到困难调动一切积极因素解决问题。}
	{\\ 具有较强的团队合作精神，性格沉稳，善于沟通，能够与团队成员和谐相处。}
		

%--------------------------------------------------------------------------------------------------
%	ARTIFICIAL FOOTER (fancy footer cannot exceed linewidth) 
%--------------------------------------------------------------------------------------------------

\null
\vspace*{\fill}
\hspace{-0.25\linewidth}\colorbox{white}{\makebox[1.5\linewidth][c]{\mystrut  \textnormal{ {Personal Page: } \textcolor{NJAUGREEN}{https://seydlitz2185.github.io} $\cdot$ {LaTeX Template: } \textcolor{NJAUGREEN}{https://github.com/jankapunkt/latexcv}}}}




%============================================================================%
%
%
%
%	DOCUMENT END
%
%
%
%============================================================================%
\end{document}
